\subsection{Further Topics in Differential Forms and Stokes Theory}

\begin{proofmethod} \label{meth:differential-forms-proof-patterns} \hlt{Differential-Forms Proof Patterns}\\
\hlt{Use When}: A proof mixes form algebra with parametrization or boundary integration.\\
Select the branch matching the obstruction and verify its hypotheses.
\begin{enumerate}[label=(\roman*)]
\setlength{\itemsep}{0pt}
\item \hlt{Signs}: Reduce wedge permutations by Proposition \ref{prop:basic-forms-standard-presentation}(i). Track determinant swaps by Theorem \ref{thm:determinant-basic-laws}(iii).
\item \hlt{Standard presentation}: Delete repeated-index terms and sort the remaining factors.\\
Existence and uniqueness of the ordered result hold by Proposition \ref{prop:basic-forms-standard-presentation}.
\item \hlt{Pullback}: Replace $dy_i$ by $dt_i$ by Definition \ref{def:pullback-of-differential-form}. Commute pullback with required operation by Theorem \ref{thm:pullback-basic-laws}.
\item \hlt{Stokes}: Identify $\lambda$ with $d\lambda$ equal to target interior form,  verify boundary orientation. Then $\int_\Psi d\lambda=\int_{\partial_{\mathrm{ch}}\Psi}\lambda$ by Theorem \ref{thm:stokes-theorem-differential-forms}.\\
For a coordinate proof, reduce by linearity to one chain and one basic summand.\\
Pull back to $Q^k$, retain only two faces associated with missing differential. Close each slice in that coordinate by Theorem \ref{thm:fundamental-theorem-calculus-riemann}.
\item \hlt{Closed versus exact}: If $\omega=d\lambda$ with $\lambda\in C^2$, then $d\omega=0$ by Theorem \ref{thm:exterior-derivative-basic-laws}(ii).\\
On a convex open set, a closed $C^1$ positive-degree $\omega$ is exact by Theorem \ref{thm:poincare-lemma-convex-rudin}(i).\\
On its $C^2$-diffeomorphic image, the same conclusion holds by Theorem \ref{thm:poincare-lemma-convex-rudin}(ii).
\item \hlt{Vector analysis}: Encode vector fields as forms by Definition \ref{def:vector-fields-as-differential-forms}.\\
Identify exterior derivatives with curl or divergence by Theorem \ref{thm:vector-analysis-forms-dictionary}.\\
Convert them to boundary identities by Theorem \ref{thm:stokes-theorem-differential-forms}. Translate the integrals by Definition \ref{def:vector-calculus-integral-notation}.
\end{enumerate}
\hlt{Key Insight}: Normal form fixes signs; pullback changes variables; Stokes moves $d$ to the boundary.
\end{proofmethod}
